\begin{frame}{Trochę historii}
\begin{columns}
\begin{column}{0.65\textwidth}
    % \textcolor{cfblue}{Dychotomia Hella-Nešetřila} \\
    %     Jeśli \(H\) jest dwudzielny albo ma pętle to \(CSP(H)\) jest w \(P\), wpp jest NP-zupełny.
    % \\\vspace{1em}
    % \onslide<2->{
    % \textcolor{cfblue}{Hipoteza (Bang-Jensen i Hell, 1990)} \\
    %     Niech \(H\) będzie spójnym grafem bez źródeł i ujść. Jeśli \(H\) ma homomorfizm do cyklu, to~problem CSP\((H)\) jest w klasie P, wpp CSP\((H)\) jest NP-zupełny.
    % }
    % \\ \vspace{1em}
    % \onslide<3->{
    %     \textcolor{cfblue}{Dychotomia CSP (Bulatov, Zhuk; 2017)} \\
    %     Dla grafu \(H\) problem CSP\((H)\) albo jest w P, albo~NP-zupełny.
    % }
    \vspace{-3em}
    \begin{figure}
        \centering
        \begin{tikzpicture}
  \node[
    rounded corners=5pt,
    inner sep=8pt,
    text width=\linewidth
  ] (theorem) {
    \textcolor{cfblue}{Dychotomia Hella-Nešetřila (1990)} \\
        Jeśli \(H\) jest dwudzielny albo ma pętle, to \(CSP(H)\) jest w~klasie \(P\), wpp \(CSP(H)\) jest NP-zupełny.

  };
\end{tikzpicture}
    \end{figure}
    \vspace{-5em}
    \onslide<2->{
        \begin{figure}
            \centering
            \begin{tikzpicture}
\only<1-3>{
  \node[
    rounded corners=5pt,
    inner sep=8pt,
    text width=\linewidth
  ] (theorem) {
    \textcolor{cfblue}{Hipoteza Bang-Jensena i Hella\\ (Kozik, Barto, Nieven 2008)}\\
    Niech \(H\) będzie spójnym grafem bez źródeł i ujść. \\
    Jeśli \(H\) ma homomorfizm do cyklu, to
    CSP\((H)\) jest w~klasie P, wpp CSP\((H)\) jest NP-zupełny.
  };
    }
  \only<4>{
    \node[
    rounded corners=5pt,
    draw=cfgrey!40,
    fill=cfgrey!40,
    inner sep=8pt,
    text width=\linewidth
  ] (theorem) {
    \textcolor{cfblue}{Hipoteza Bang-Jensena i Hella\\ (Kozik, Barto, Nieven 2008)}\\
    Niech \(H\) będzie spójnym grafem bez źródeł i ujść. \\
    Jeśli \(H\) ma homomorfizm do cyklu, to problem
    CSP\((H)\) jest w~klasie P, wpp CSP\((H)\) jest NP-zupełny.
  };
  }
\end{tikzpicture}
        \end{figure}
    }
    \vspace{-5em}
    \onslide<3->{
    \begin{figure}
        \centering
        \begin{tikzpicture}
  \node[
    rounded corners=5pt,
    inner sep=8pt,
    text width=\linewidth
  ] (theorem) {
    \textcolor{cfblue}{Dychotomia CSP (Bulatov, Zhuk; 2017)} \\
    Dla grafu \(H\) problem CSP\((H)\) albo jest w P, albo~jest NP-zupełny.
    }
\end{tikzpicture}
    \end{figure}
    }
\end{column}
\begin{column}{0.35\textwidth}
    \begin{figure}[H]
        \centering
        \begin{tikzpicture}[scale=1, line width=1.2pt]

% Jeśli P \cap NP-hard = \emptyset to P\neq NP
% ----------- NP ------------
\draw (0,0) ellipse (1.75 and 1.45);
\node at (1.2,0) {\large NP};

% ------- NP-complete -------
\draw (0,0.75) ellipse (0.8 and 0.55);
\node at (1.4,1.75) {\large NP-zupełne};
\draw (1.35,1.45) -- (0.5,0.8);
% \node at (-0.35,-0.8) {\large P};

% ------------ P ------------
\draw (0,-0.75) ellipse (0.8 and 0.55);
\node at (0.4,-0.8) {\large P};

% ----------- CSP -----------
\begin{scope}[line width=1.6pt, rounded corners=8pt, cfblue]
\only<1-2>{
    \draw (-0.55,-1) rectangle (0,1);
    \node at (-1.1,0) {\Large CSP};
}

\only<3->{
    \draw (-0.55,0.4) rectangle (0,1);
    \draw (-0.55,-1) rectangle (0,-0.4);
    \node at (-1.1,0) {\Large CSP};
}
\end{scope}

\end{tikzpicture} 
    \end{figure}
    \vspace{1em}
\end{column}
\end{columns}

\end{frame}